\documentclass[openany,fontset=fandol]{ctexbook}
\usepackage{amsmath,amssymb}
\usepackage[unicode,hidelinks]{hyperref}
\newcommand{\note}[1]{\footnote{译注：#1}}
\begin{document}
\title{代数几何基础第五卷预备笔记\\第12节：奇异关联簇与交维数过滤}
\author{据用户上传英译编辑稿译出}
\date{2026年11月16日}
\maketitle
\chapter*{范围与校订状态}
本批完整译出跨底本第46—47页的一长段，自“Let us again assume that”起，止于“a `geometric' paragraph devoted to grassmanians”。为便于阅读，按论证层次分为四段。下一接点是第47页“If we also define $V^{(k)}$”。

本段含原作者的试探性判断、难辨标记及未定位的补篇引用，均保留；不把尚待验证的过滤与光滑性陈述改写成已证明定理。本稿独立保存，未合入主稿。此前待核事项不因本批而消除。
\chapter*{§12 将上述结果推广到线性截面（续）}
再假设$f\colon X\to P$未分歧，且$X$在$S$上光滑，各分支的维数至少为$m$。于是$X^{(m)}=V^0$；我们在其中区分出子预概形
\[
X^{(m)\mathrm{sing}}=V^1,
\]
它由$\phi^m$相对于$\operatorname{Grass}^m$的奇异零点组成。从几何上看，它也由这样的点对$(x,\xi)$组成：线性子簇$L_\xi$与$X$在$x$处的切空间（视为$P$的线性子簇）的交维数超出预期，也就是说，这两个空间不能张成整个$P$。\note{底本此处确实使用$V^0$、$V^1$和$\phi^m$，而第47页转用$V'$、$V''$。正文保留原记号差别。$X^{(m)\mathrm{sing}}$在此指相对于参数空间的奇异零点子预概形，不是对总空间绝对奇异点集的断言；切空间叙述是在相应几何纤维中理解的。}

与$m=1$时的情形相反，对于任意$m$，态射
\[
X^{(m)\mathrm{sing}}\longrightarrow X
\]
一般不光滑。这是因为，经过$x$〔字迹不清，请询问A.G.〕、并与给定线性子簇$T\ni x$相交而交维数超出预期的那些$L$所组成的簇，一般并不在$k$上光滑：这个簇〔字迹不清，请询问A.G.〕，只有由如下$\xi$组成的光滑子簇〔此处句意残缺〕：$T\cap L_\xi$的维数恰好比“正常”维数$n-m$大一，其中$n=\dim T$、$m=\operatorname{codim}L$。〔V；法文第71页。〕\note{底本在两处“illegible, ask AG”之外，还印有无法确定词义及句法的“only the loneve of the subvariety smooth”。因此这里不擅自补成“闭包”“光滑部分”或其他确定关系；可辨的交维数描述仍完整保留。前面的“V [page 71 Fr]”也按底本保留，未把V猜定为交叉引用或新对象。}

\newpage
如果没有弄错，第16段补篇中引入的集合$V''$（它包含在相对超奇异点集中），无非就是由满足
\[
\dim(T_x\cap L_\xi)\geq n-m+2
\]
的点对$(x,L_\xi)$组成的集合。因此，$V'\setminus V''$在$S$上光滑；如果没有弄错，它也恰好是$V'$在$X$上光滑的那些点组成的集合。\note{原作者在本句两次使用“barring an error”。这里据图像核为撇号$V'$、$V''$，不是$V^0$和$V^{00}$；不将文本抽取的撇号误读沿入正文。第46页的$V^1$与这里$V'$的对应虽由语境提示，底本未明确重定义，本批不悄然统一。}

（要验证这一点，需要按照交的维数研究格拉斯曼概形的过滤，其中$L$变动而$T$固定。如果没有弄错，当我们不仅在集合意义下，而且还借助第16段补篇第16页的引理，在概形意义下定义这一过滤时，过滤的下一层〔71〕恰好由前一层的非光滑点组成。〔此处原译者在“层”后注：法文？分层中的层？？？〕\footnote{原脚注61：询问A.G.。}\footnote{原脚注62：A.G.，请帮助定位这一引用。}因此，这项研究将构成专门讨论格拉斯曼簇的一个“几何性”段落中的一小节。）\note{底本使用“notch”，并附“(*) [Fr] (stratum???)”；“层”在这里仅为保留过滤语境的译法，不判定它指闭过滤项还是相邻项之差。原脚注61、62与无法定位的补篇引理均保留；没有据此补写过滤的概形结构或其非光滑点定理。}
\end{document}
